% Source: https://artofproblemsolving.com/community/c3h2265593_number_theory_problem
% Stable identity: aops-topic-2265593-q01
Clearing denominators, $3n+3m=mn$, so $mn-3m-3n+9=9$, that is,
\[(m-3)(n-3)=9.\]
Since $\dfrac{1}{m}<\dfrac{1}{3}$ is needed (as $\dfrac{1}{n}>0$), we have $m>3$, so both factors are positive with $m-3<n-3$. The only possibility is $m-3=1$, $n-3=9$: $m=4$, $n=12$. (The factorisation $3\cdot3$ would give $m=n$.) Then
\[\dfrac{1}{m}-\dfrac{1}{n}=\dfrac{1}{4}-\dfrac{1}{12}=\dfrac{2}{12}=\boxed{\textbf{(C)}\ \dfrac{1}{6}}.\]
